\par
\SourceLeaf{368}%
Ergeben drei aufeinanderfolgende Skalen
\[
M,0,0,0
\]
und sind die betreffenden Drehungspunkte $P,P',P''$,\\
\hspace*{5em}die Winkel $\lambda,\lambda',\lambda''$,\\
so sind $\tfrac12\lambda,\tfrac12\lambda',\tfrac12\lambda''$ die Winkel des sphärischen Dreiecks zwischen $P,P',P''$.

\Needspace{14\baselineskip}
\begin{center}[II.]\end{center}
\begin{center}\bfseries MUTATIONEN DES RAUMES VON DREI DIMENSIONEN.\end{center}
\begin{center}1.\end{center}

Seit langer Zeit verwende ich das Verfahren, jede Mutation eines Raums (darunter verstehe ich seine Verschiebung in einem anderen Raum verbunden mit einer allgemeinen Vergrößerung oder Verkleinerung des ersteren unter Wahrung der Ähnlichkeit) durch vier Größen
\[
a,b,c,d
\]
auszudrücken, deren Gesamtheit ich die Mutationsskala nenne. Zugleich wird angenommen, dass ein Punkt des beweglichen Raums in dem anderen, der als fest oder absolut betrachtet wird, stets unbewegt bleibt. Dieser Punkt mag als Nullpunkt bezeichnet werden.

\Needspace{9\baselineskip}
\begin{center}2.\end{center}

Diese vier Größen haben dabei folgende Bedeutung. Im festen Raum nehme man drei von $0$ ausgehende Geraden an, die miteinander rechte Winkel bilden, $01,02,03$. Man setze
\[
\BareRoot{(bb+cc+dd)}=\rho,
\]
und lasse
\[
\frac b\rho,\ \frac c\rho,\ \frac d\rho
\]
